\subsection{Systematic Underconfidence as a Structural Property}
\label{sec:discussion:underconfidence}

The dominant finding of this study is that ALeRCE's light curve 
classifier exhibits systematic underconfidence rather than the 
overconfidence commonly reported for deep learning classifiers 
\citep{Guo2017, Minderer2021}. This distinction matters for 
operational astronomy. Overconfident classifiers produce dangerous 
silent failures — high-confidence wrong predictions that are acted 
upon without human review. Underconfident classifiers produce 
unnecessary deferrals — objects the classifier effectively 
discriminates but assigns low confidence, causing them to be 
flagged for spectroscopic follow-up that is not scientifically 
required.

In the resource-constrained regime of LSST-era astronomy, where 
spectroscopic follow-up capacity is approximately $10^5$ 
observations per year against $10^7$ nightly alerts, unnecessary 
deferrals carry a direct opportunity cost. Every spectrum allocated 
to an object the classifier could have confidently classified 
correctly is a spectrum not allocated to a genuinely ambiguous 
object or a rare high-priority target. Quantifying and correcting 
underconfidence is therefore not merely a technical calibration 
exercise but a resource allocation problem.

\subsection{Why Random Forest Suppresses Confidence}
\label{sec:discussion:rf}

The underconfidence we observe is mechanistically expected for 
Random Forest classifiers in class-imbalanced settings. RF 
probability estimates are computed as the fraction of trees 
voting for each class. In a forest trained on an imbalanced 
dataset — SNIa comprises 47\% of our sample — the decision 
boundaries are dominated by the majority class, and minority 
class probabilities are systematically suppressed \citep{Niculescu-Mizil2005}. 
This suppression is strongest for classes the classifier 
discriminates most reliably, because reliable discrimination 
concentrates vote fractions near the boundary between correct 
and incorrect, rather than near 1.0.

This explains the inverse relationship we observe between accuracy 
and ECE across classes. SNIa and SNIbc are classified with 
accuracies above 0.83, yet their ECEs exceed 0.36. SNII, 
classified at only 0.495 accuracy, has ECE of 0.103. The 
classifier is most miscalibrated precisely where it has 
the most genuine discriminative signal.

The practical implication is that RF-based astronomical classifiers 
should not be used for confidence-based deferral decisions without 
post-hoc calibration. The raw probability outputs do not represent 
reliable confidence estimates, even when classification accuracy 
is high.

\subsection{Temperature Scaling: Efficacy and Limits}
\label{sec:discussion:tempscaling}

Temperature scaling with $T = 0.359$ reduces ALeRCE's ECE from 
$0.259$ to $0.031$ — an 88\% improvement with a single scalar 
parameter requiring no model retraining. This result is 
encouraging for operational deployment: existing broker 
infrastructure could be augmented with a post-hoc calibration 
layer at negligible computational cost.

However, the per-class analysis reveals an important limitation. 
Global temperature scaling worsens SNII calibration from $0.099$ 
to $0.189$, because SNII was already well-calibrated before 
scaling. A single $T$ cannot simultaneously optimise calibration 
for classes with heterogeneous miscalibration structure. 

Two approaches could address this. First, class-specific 
temperature parameters — fitting a separate $T_k$ for each 
class $k$ — would allow independent correction of each 
class's calibration curve. Second, more expressive post-hoc 
methods such as isotonic regression \citep{Zadrozny2002} or 
matrix scaling \citep{Kull2019} could model class interactions 
that scalar temperature cannot capture. We leave systematic 
comparison of these methods to future work, noting that the 
simple temperature scaling result already achieves practically 
useful calibration for the majority of classes.

\subsection{Fink's Structural Miscalibration}
\label{sec:discussion:fink}

The failure of temperature scaling for Fink's RF classifier — 
evidenced by the optimiser hitting the upper bound at $T \approx 10$ 
— indicates a qualitatively different problem from ALeRCE's 
underconfidence. When no finite temperature produces a well-calibrated 
output, the score distribution lacks the monotonic relationship 
between predicted probability and empirical accuracy that 
temperature scaling requires.

Inspection of the score distribution confirms this: 
\texttt{rf\_snia\_vs\_nonia} assigns scores below $0.05$ to 
over 90\% of objects, including the majority of true SNIa. 
The classifier is not underconfident in the same sense as 
ALeRCE — it is effectively non-discriminative for SNIa 
identification on this dataset. This may reflect a domain 
shift between Fink's training data and the BTS sample, or 
a fundamental limitation of the binary RF formulation for 
this classification problem.

SuperNNova retains weak discriminative structure — its 
reliability curve rises slightly with confidence — but 
insufficient for effective calibration. The deep learning 
architecture of SuperNNova provides some advantage over 
the RF approach, consistent with broader trends in the 
literature \citep{Moller2019, Villar2019}.

\subsection{Implications for Deferral-Aware Classification}
\label{sec:discussion:deferral}

This calibration study motivates the deferral-aware 
classification framework that constitutes the broader 
research programme of which this paper is the first component. 
Learning-to-defer frameworks \citep{Madras2018, Mozannar2020} 
require reliable confidence estimates as input to the deferral 
policy: an object should be deferred when classifier uncertainty 
genuinely reflects the difficulty of the classification, not 
when a well-discriminating classifier mechanically suppresses 
its own signal.

Our results suggest that ALeRCE, post-calibration, provides 
confidence estimates of sufficient quality to support deferral 
decisions ($\mathrm{ECE} = 0.031$). Fink's classifiers, in 
their current form, do not. This finding directly informs 
broker selection for downstream deferral framework development.

The TDE taxonomy gap identified in Section~\ref{sec:results:tde} 
represents a separate but related challenge. Deferral frameworks 
require that all scientifically relevant classes appear in the 
classifier's output space. A broker that cannot classify TDE 
by design will never defer TDE candidates for spectroscopic 
follow-up — it will instead misclassify them as the nearest 
available class with high confidence, producing exactly the 
silent failure mode that deferral is designed to prevent.

\subsection{Limitations}
\label{sec:discussion:limitations}

Several limitations of this study should be acknowledged. 
First, we evaluate calibration on a single broker snapshot; 
ALeRCE and Fink update their classifiers periodically, and 
our results reflect classifier versions active during our 
data collection period. Second, the BTS magnitude limit 
restricts our sample to bright, low-redshift transients; 
calibration properties may differ for fainter populations 
that will dominate LSST alert streams. Third, our temperature 
scaling analysis is conducted on the evaluation set itself, 
providing an optimistic estimate of post-hoc calibration 
performance; deployment would require fitting $T$ on a 
held-out validation set.

Finally, we note that ECE as defined in 
Equation~\ref{eq:ece} is sensitive to bin width and the 
distribution of samples across bins \citep{Nixon2019}. 
Our bootstrap confidence intervals partially address this 
by quantifying sampling uncertainty, but do not capture 
binning sensitivity. Adaptive binning methods \citep{Roelofs2022} 
may provide more robust ECE estimates for highly imbalanced 
confidence distributions such as those observed for Fink's 
RF classifier.